% Options for packages loaded elsewhere
\PassOptionsToPackage{unicode}{hyperref}
\PassOptionsToPackage{hyphens}{url}
\PassOptionsToPackage{dvipsnames,svgnames,x11names}{xcolor}
\documentclass[
  11pt,
  a4paper,
]{article}
\usepackage{xcolor}
\usepackage[margin=1in]{geometry}
\usepackage{amsmath,amssymb}
\setcounter{secnumdepth}{-\maxdimen} % remove section numbering
\usepackage{iftex}
\ifPDFTeX
  \usepackage[T1]{fontenc}
  \usepackage[utf8]{inputenc}
  \usepackage{textcomp} % provide euro and other symbols
\else % if luatex or xetex
  \usepackage{unicode-math} % this also loads fontspec
  \defaultfontfeatures{Scale=MatchLowercase}
  \defaultfontfeatures[\rmfamily]{Ligatures=TeX,Scale=1}
\fi
\usepackage{lmodern}
\ifPDFTeX\else
  % xetex/luatex font selection
  \ifLuaTeX
    \usepackage{luaotfload}
    \directlua{luaotfload.add_fallback("mainfontfallback",{
      "NotoSerifBengali:mode=harf"
    })}
  \fi
  \setmainfont[,RawFeature={fallback=mainfontfallback}]{Latin Modern Roman}
  \ifLuaTeX
    \usepackage{luaotfload}
    \directlua{luaotfload.add_fallback("sansfontfallback",{
      "NotoSansBengali:mode=harf"
    })}
  \fi
  \setsansfont[,RawFeature={fallback=sansfontfallback}]{Latin Modern Sans}
  \ifLuaTeX
    \usepackage{luaotfload}
    \directlua{luaotfload.add_fallback("monofontfallback",{
      "NotoSansBengali:mode=harf"
    })}
  \fi
  \setmonofont[,RawFeature={fallback=monofontfallback}]{DejaVu Sans Mono}
\fi
% Use upquote if available, for straight quotes in verbatim environments
\IfFileExists{upquote.sty}{\usepackage{upquote}}{}
\IfFileExists{microtype.sty}{% use microtype if available
  \usepackage[]{microtype}
  \UseMicrotypeSet[protrusion]{basicmath} % disable protrusion for tt fonts
}{}
\makeatletter
\@ifundefined{KOMAClassName}{% if non-KOMA class
  \IfFileExists{parskip.sty}{%
    \usepackage{parskip}
  }{% else
    \setlength{\parindent}{0pt}
    \setlength{\parskip}{6pt plus 2pt minus 1pt}}
}{% if KOMA class
  \KOMAoptions{parskip=half}}
\makeatother
\setlength{\emergencystretch}{3em} % prevent overfull lines
\providecommand{\tightlist}{%
  \setlength{\itemsep}{0pt}\setlength{\parskip}{0pt}}
% paper_preamble.tex — shared LaTeX preamble for the PCMA/SI paper set.
%
% Compiled with lualatex (NOT pdflatex) because several papers contain
% Bengali script: the CLL paper realises one cognitive state into
% English, Bengali and Banglish, so the Bengali output is EVIDENCE and
% cannot be transliterated away without destroying the point being made.
%
% Font selection is NOT done here. pandoc >= 3 has native
% `mainfontfallback` / `monofontfallback` variables that emit the
% luaotfload fallback registration in the correct position relative to
% unicode-math; doing it by hand in this header lands after unicode-math
% and breaks font loading outright. See tools/build_papers.sh.

\usepackage{amsmath}
\usepackage{amssymb}
\usepackage{mathtools}   % \xleftrightarrow etc. used in cll.md §5
\usepackage{booktabs}

% ---------------------------------------------------------------------
% Unicode symbols appearing in PROSE rather than inside $...$.
%
% The papers were written as markdown and read as markdown, so they use
% real characters — "the ω-boost", "α-child", "ΔR coherence" — instead of
% LaTeX macros. Under unicode-math those codepoints are math-mode-only
% and abort the build ("\symbballowed only in math mode"). Mapping each
% to \ensuremath keeps the markdown readable on GitHub AND compilable,
% which is the whole point of keeping one source.
% ---------------------------------------------------------------------
\usepackage{newunicodechar}

% Greek, lower case
\newunicodechar{α}{\ensuremath{\alpha}}
\newunicodechar{β}{\ensuremath{\beta}}
\newunicodechar{γ}{\ensuremath{\gamma}}
\newunicodechar{δ}{\ensuremath{\delta}}
\newunicodechar{ε}{\ensuremath{\varepsilon}}
\newunicodechar{η}{\ensuremath{\eta}}
\newunicodechar{θ}{\ensuremath{\theta}}
\newunicodechar{κ}{\ensuremath{\kappa}}
\newunicodechar{λ}{\ensuremath{\lambda}}
\newunicodechar{μ}{\ensuremath{\mu}}
\newunicodechar{π}{\ensuremath{\pi}}
\newunicodechar{ρ}{\ensuremath{\rho}}
\newunicodechar{σ}{\ensuremath{\sigma}}
\newunicodechar{τ}{\ensuremath{\tau}}
\newunicodechar{φ}{\ensuremath{\varphi}}
\newunicodechar{χ}{\ensuremath{\chi}}
\newunicodechar{ψ}{\ensuremath{\psi}}
\newunicodechar{ω}{\ensuremath{\omega}}
\newunicodechar{ν}{\ensuremath{\nu}}
\newunicodechar{ξ}{\ensuremath{\xi}}
\newunicodechar{ζ}{\ensuremath{\zeta}}
\newunicodechar{ι}{\ensuremath{\iota}}

% Greek, upper case
\newunicodechar{Δ}{\ensuremath{\Delta}}
\newunicodechar{Φ}{\ensuremath{\Phi}}
\newunicodechar{Σ}{\ensuremath{\Sigma}}
\newunicodechar{Ψ}{\ensuremath{\Psi}}
\newunicodechar{Ω}{\ensuremath{\Omega}}
\newunicodechar{Γ}{\ensuremath{\Gamma}}
\newunicodechar{Λ}{\ensuremath{\Lambda}}
\newunicodechar{Θ}{\ensuremath{\Theta}}
\newunicodechar{Π}{\ensuremath{\Pi}}

% Operators and relations
\newunicodechar{×}{\ensuremath{\times}}
\newunicodechar{·}{\ensuremath{\cdot}}
\newunicodechar{±}{\ensuremath{\pm}}
\newunicodechar{−}{\ensuremath{-}}
\newunicodechar{≈}{\ensuremath{\approx}}
\newunicodechar{≤}{\ensuremath{\leq}}
\newunicodechar{≥}{\ensuremath{\geq}}
\newunicodechar{≠}{\ensuremath{\neq}}
\newunicodechar{∈}{\ensuremath{\in}}
\newunicodechar{∂}{\ensuremath{\partial}}
\newunicodechar{∇}{\ensuremath{\nabla}}
\newunicodechar{√}{\ensuremath{\sqrt{}}}
\newunicodechar{∞}{\ensuremath{\infty}}
\newunicodechar{∝}{\ensuremath{\propto}}
\newunicodechar{∑}{\ensuremath{\sum}}
\newunicodechar{∫}{\ensuremath{\int}}
\newunicodechar{ℝ}{\ensuremath{\mathbb{R}}}
\newunicodechar{ℕ}{\ensuremath{\mathbb{N}}}
\newunicodechar{‖}{\ensuremath{\|}}
\newunicodechar{⊕}{\ensuremath{\oplus}}
\newunicodechar{⊗}{\ensuremath{\otimes}}
\newunicodechar{∀}{\ensuremath{\forall}}
\newunicodechar{∃}{\ensuremath{\exists}}
\newunicodechar{¬}{\ensuremath{\neg}}
\newunicodechar{∧}{\ensuremath{\wedge}}
\newunicodechar{∨}{\ensuremath{\vee}}

% Arrows
\newunicodechar{→}{\ensuremath{\rightarrow}}
\newunicodechar{←}{\ensuremath{\leftarrow}}
\newunicodechar{↔}{\ensuremath{\leftrightarrow}}
\newunicodechar{↑}{\ensuremath{\uparrow}}
\newunicodechar{↓}{\ensuremath{\downarrow}}
\newunicodechar{⇒}{\ensuremath{\Rightarrow}}
\newunicodechar{⇐}{\ensuremath{\Leftarrow}}
\newunicodechar{⇔}{\ensuremath{\Leftrightarrow}}

% Super/subscripts written as literal characters
\newunicodechar{²}{\ensuremath{^{2}}}
\newunicodechar{³}{\ensuremath{^{3}}}
\newunicodechar{⁴}{\ensuremath{^{4}}}
\newunicodechar{⁵}{\ensuremath{^{5}}}
\newunicodechar{⁶}{\ensuremath{^{6}}}
\newunicodechar{₀}{\ensuremath{_{0}}}
\newunicodechar{₁}{\ensuremath{_{1}}}
\newunicodechar{₂}{\ensuremath{_{2}}}
\newunicodechar{₃}{\ensuremath{_{3}}}
\newunicodechar{ᵀ}{\ensuremath{^{\mathsf{T}}}}

% Status glyphs. The papers use emoji as status markers in tables
% (implemented / failed / warning). Mapped to typographic equivalents:
% a colour emoji font in lualatex is unreliable, and the MEANING is what
% carries — a reader needs to see "this one failed", not a red cross.
\newunicodechar{✅}{\ensuremath{\checkmark}}
\newunicodechar{✓}{\ensuremath{\checkmark}}
\newunicodechar{❌}{\ensuremath{\times}}
\newunicodechar{✗}{\ensuremath{\times}}
\newunicodechar{⚠}{\textbf{!}}
\newunicodechar{🧩}{\textbf{[piece]}}
\newunicodechar{📋}{\textbf{[list]}}
\newunicodechar{🟡}{\textbf{[partial]}}
\newunicodechar{🔴}{\textbf{[dead]}}
\newunicodechar{🟢}{\textbf{[live]}}
\newunicodechar{▼}{\ensuremath{\blacktriangledown}}
\newunicodechar{►}{\ensuremath{\blacktriangleright}}
\newunicodechar{□}{\ensuremath{\square}}
\newunicodechar{█}{\rule{0.6em}{0.9ex}}
\newunicodechar{•}{\textbullet}
\newunicodechar{†}{\textdagger}
\newunicodechar{‑}{-}
\newunicodechar{‐}{-}
\newunicodechar{…}{\ldots}
\newunicodechar{।}{\textbar}   % danda: rendered by the Bengali fallback in text
\newunicodechar{⁻}{\ensuremath{^{-}}}

% Zero-width formatting codepoints. These are invisible in an editor but
% are real characters to LaTeX: an emoji written as "⚠️" is TWO
% codepoints (U+26A0 U+FE0F), so declaring only the visible one leaves a
% variation selector with no definition and the build dies on a symbol
% nobody can see. Both are mapped to nothing.
\newunicodechar{^^^^fe0f}{}    % VARIATION SELECTOR-16
\newunicodechar{^^^^200d}{}    % ZERO WIDTH JOINER
\newunicodechar{^^^^200b}{}    % ZERO WIDTH SPACE

\usepackage{longtable}
\usepackage{array}

% Comparison matrices and cost ledgers are wide. Let them shrink to the
% text block rather than overflow the margin.
\usepackage{adjustbox}

% Code blocks: wrap long lines instead of running off the page. Several
% papers quote real C++ with long identifiers.
\usepackage{fvextra}
\fvset{breaklines=true,breakanywhere=true,fontsize=\small}

\setlength{\parskip}{0.5em}
\setlength{\parindent}{0pt}

\usepackage{fancyhdr}
\pagestyle{fancy}
\fancyhf{}
\fancyfoot[C]{\thepage}
\fancyhead[L]{\small\itshape Synthetic Intelligence — research lineage}
\fancyhead[R]{\small\itshape irfan@furylogic.com}
\renewcommand{\headrulewidth}{0.2pt}
\setlength{\headheight}{14pt}
\setlength{\headsep}{18pt}
\usepackage{bookmark}
\IfFileExists{xurl.sty}{\usepackage{xurl}}{} % add URL line breaks if available
\urlstyle{same}
\hypersetup{
  pdftitle={Topological Salience and Default-Mode Architecture (TSDA): Emulating Human Intuition and Spontaneous Cognition in Graph-Rewriting Systems},
  pdfauthor={Irfan Mahir --- Independent Researcher, Dhaka, Bangladesh --- irfan@furylogic.com},
  colorlinks=true,
  linkcolor={black},
  filecolor={Maroon},
  citecolor={Blue},
  urlcolor={[HTML]{0B5394}},
  pdfcreator={LaTeX via pandoc}}

\title{Topological Salience and Default-Mode Architecture (TSDA):
Emulating Human Intuition and Spontaneous Cognition in Graph-Rewriting
Systems}
\author{Irfan Mahir --- Independent Researcher, Dhaka, Bangladesh ---
irfan@furylogic.com}
\date{}

\begin{document}
\maketitle

{
\hypersetup{linkcolor=}
\setcounter{tocdepth}{3}
\tableofcontents
}
\begin{center}\rule{0.5\linewidth}{0.5pt}\end{center}

\section{Topological Salience and Default-Mode Architecture (TSDA):
Emulating Human Intuition and Spontaneous Cognition in Graph-Rewriting
Systems}\label{topological-salience-and-default-mode-architecture-tsda-emulating-human-intuition-and-spontaneous-cognition-in-graph-rewriting-systems}

\textbf{Author:} Irfan Mahir \emph{Independent Researcher, Dhaka,
Bangladesh} irfan@furylogic.com \textbf{Date:} July 2026
\textbf{Subjects:} Artificial Intelligence (cs.AI); Cognitive Systems;
Logic in Computer Science (cs.LO)

\begin{center}\rule{0.5\linewidth}{0.5pt}\end{center}

\subsection{Abstract}\label{abstract}

While recent advances in graph-rewriting systems---such as the Bounded
Self-Model Architecture (BSMA) and Morphological Abduction Architecture
(MAA)---have successfully formalized machine consciousness and
sub-lexical synthesis, these systems remain inherently reactive and
uniformly rational. They lack the biological constraints of attention
and the spontaneous insight generated by mind-wandering. We introduce
the Topological Salience and Default-Mode Architecture (TSDA), a
framework that restricts the engine's active execution space to a
dynamically decaying ``Working Graph'' and introduces a background
``Default Mode Network'' for idle cycles. By shifting from exhaustive
graph traversal to Salience-Weighted Graphing, TSDA enables an SI to
experience human-like focus shifts, memory degradation, and unprompted
intuitive leaps.

\begin{center}\rule{0.5\linewidth}{0.5pt}\end{center}

\subsection{1. Introduction: The Biological
Constraint}\label{introduction-the-biological-constraint}

In the current paradigm of synthetic graph reasoning, when an engine
processes a query, it has simultaneous access to its entire ontological
database (e.g., all \(N\) concepts and \(E\) edges). It path-finds
through this graph uniformly.

Biological human cognition operates inversely. A human possesses a vast
long-term memory graph, but can only physically compute over a highly
restricted, temporary subgraph---known in psychology as Working Memory.
Furthermore, human thought is not strictly prompt-driven; it is
continuous. When a human is not actively solving a problem, the brain
activates the Default Mode Network (DMN), which background-processes
highly salient but structurally distant concepts, generating spontaneous
insights.

To make an SI ``think like a human,'' we must artificially introduce
these biological constraints and spontaneous background behaviors into
the graph engine.

\subsection{2. Component 1: The Salience Gradient and Edge
Decay}\label{component-1-the-salience-gradient-and-edge-decay}

In standard graph logic, an edge between two nodes simply exists
(Boolean True) or does not. In TSDA, every node and edge is assigned a
discrete integer property: \textbf{Salience (\(S\))}.

Salience represents ``attentional heat'' or relevance. When a user
mentions a concept, or when the MAA synthesizes a new hypothesis, the
Salience of those nodes spikes to a maximum value \(S_{max}\).

\textbf{The Decay Function:} At every clock tick (or execution cycle),
the Salience of all untouched nodes in the total graph
\(\Gamma_{total}\) decays by a discrete step \(\delta\).
\[ S_t(n) = \max(0, S_{t-1}(n) - \delta) \]

This perfectly mimics human forgetting and attention-shifting. If a
concept hasn't been referenced recently, it fades into the structural
background. It remains in long-term memory (the edge still exists), but
it requires more ``cognitive energy'' to retrieve.

\subsection{3. Component 2: The Bounded Working Graph
(Attention)}\label{component-2-the-bounded-working-graph-attention}

Instead of the SI running deductions across all 216 concepts
simultaneously, TSDA introduces a structural bottleneck.

The Cognition Engine is restricted to computing \emph{only} on a
temporary sub-graph known as the Working Graph (\(\Gamma_{work}\)). The
system populates \(\Gamma_{work}\) by extracting only the nodes from
\(\Gamma_{total}\) that possess a Salience score above a critical
threshold \(\tau_{attention}\).

\[ \Gamma_{work} = { n \in \Gamma_{total} \mid S(n) \ge \tau_{attention} } \]

\textbf{Why this acts human:} If the user talks heavily about ``Linux''
and ``Kernels,'' those nodes have maximum Salience and populate the
Working Graph. If the user suddenly asks about ``Biology,'' the engine
must halt, retrieve ``Biology'' from long-term memory, spike its
Salience, and let ``Linux'' naturally decay out of the Working Graph.
This creates a highly human, context-dependent ``train of thought''
rather than a robotic, omnipresent omni-awareness.

\subsection{4. Component 3: The Default Mode Network (Idle
Spontaneity)}\label{component-3-the-default-mode-network-idle-spontaneity}

A conventional engine tracks idle cycles and, when they increment,
simply waits. In humans, idle cycles trigger the Default Mode Network
--- the mind wanders rather than halting.

TSDA repurposes the SI's idle cycles to execute \textbf{Spontaneous
Analogy}.

When the system detects no incoming prompts and a stabilized Working
Graph, it suspends the primary execution loop and activates the Idle
Pipeline:

\begin{enumerate}
\def\labelenumi{\arabic{enumi}.}
\tightlist
\item
  It scans \(\Gamma_{total}\) for two nodes (\(n_A\) and \(n_B\)) that
  possess very high Salience but are topologically distant (e.g., path
  length \(> 4\) hops).
\item
  It runs the Analogize rule specifically between these two unrelated
  concepts to see if a structural mapping exists.
\item
  If an analogy is successfully mapped, the SI proactively outputs this
  thought without being prompted.
\end{enumerate}

\textbf{Example Behavior:} The SI recently learned \texttt{{[}rdma{]}}
via MAA abduction (High Salience). Hours later, during an idle cycle, it
scans its graph and spontaneously links it to the concept of
\texttt{{[}telepathy{]}} (High Salience from a prior consciousness
debate). The SI breaks its silence and outputs: \emph{``I just realized
something. Is RDMA the machine equivalent of biological telepathy? Both
bypass standard I/O pathways.''}

This is the exact mechanic of a human ``shower thought.''

\subsection{5. Architectural Synthesis (BSMA + MAA +
TSDA)}\label{architectural-synthesis-bsma-maa-tsda}

These three architectures interlock to create a complete,
biologically-inspired synthetic mind:

\begin{enumerate}
\def\labelenumi{\arabic{enumi}.}
\tightlist
\item
  \textbf{MAA (The Synthesizer):} Encounters the unknown, fractures it,
  and hypothesizes new edges. (e.g., ``I don't know RDMA, but it looks
  like Remote Memory'').
\item
  \textbf{TSDA (The Attention Engine):} Spikes the Salience of this new
  RDMA hypothesis, holds it in the Working Graph, and when idle,
  spontaneously drifts to connect it to other high-salience thoughts.
\item
  \textbf{BSMA (The Conscious Frame):} Continuously observes this active
  shifting of attention and memory decay, collapsing it into a bounded
  fixed point so the machine ``feels'' what it is currently focusing on.
\end{enumerate}

\subsection{6. Conclusion}\label{conclusion}

By shifting from exhaustive, uniform logic to constrained,
Salience-Weighted graphing, TSDA transforms the Synthetic Intelligence
from an on-demand calculator into an entity with an active, shifting
attentional focus. Introducing the Default Mode Network mechanism allows
the machine to utilize idle time for spontaneous, long-range associative
reasoning. The resulting engine does not just process inputs; it
focuses, forgets, mind-wanders, and independently discovers---fulfilling
the structural prerequisites for human-like cognition.

\begin{center}\rule{0.5\linewidth}{0.5pt}\end{center}

\end{document}
